\documentclass[11pt,a4paper]{article}
\usepackage[utf8]{inputenc}
\usepackage[T1]{fontenc}
\usepackage[portuguese]{babel}
\usepackage{geometry}
\usepackage{helvet}
\usepackage{enumitem}
\usepackage{titlesec}
\usepackage[table]{xcolor}
\usepackage{hyperref}
\usepackage{tabularx}
\usepackage{booktabs}
\usepackage{fancyhdr}
\usepackage[most]{tcolorbox}
\usepackage{multicol}
\usepackage{ragged2e}
\usepackage{graphicx}
\usepackage{lastpage}

\geometry{margin=2.2cm, top=2.8cm, bottom=2.5cm}
\setlength{\headheight}{15pt}
\renewcommand{\familydefault}{\sfdefault}

% ── Paleta de cores ──
\definecolor{uefsdark}{RGB}{0,51,102}
\definecolor{uefsblue}{RGB}{25,84,150}
\definecolor{uefslight}{RGB}{232,241,250}
\definecolor{uefsgray}{RGB}{110,110,110}
\definecolor{uefsaccent}{RGB}{180,60,40}

% ── Seções ──
\titleformat{\section}
  {\normalfont\large\bfseries\color{uefsdark}\scshape}{}{0em}{}%%
  [\vspace{2pt}{\color{uefsdark}\titlerule}]
\titlespacing{\section}{0pt}{20pt}{8pt}

\titleformat{\subsection}
  {\normalfont\normalsize\bfseries\color{uefsblue}}{}{0em}{}
\titlespacing{\subsection}{0pt}{12pt}{4pt}

% ── Listas ──
\setlist[itemize]{leftmargin=1.5em,topsep=4pt,itemsep=2pt,parsep=0pt}
\setlist[enumerate]{leftmargin=2em,topsep=4pt,itemsep=4pt,parsep=0pt}

% ── tcolorbox ──
\tcbset{
  uefsbox/.style={
    colback=uefslight, colframe=uefsdark,
    boxrule=0.4pt, arc=3pt,
    left=8pt, right=8pt, top=6pt, bottom=6pt,
    fontupper=\small
  },
  ementabox/.style={
    colback=white, colframe=uefsdark,
    boxrule=0.8pt, arc=0pt,
    left=10pt, right=10pt, top=8pt, bottom=8pt,
    title={\scshape\bfseries Ementa},
    fonttitle=\large\color{white},
    coltitle=white, colbacktitle=uefsdark,
    attach boxed title to top left={yshift=-2mm,xshift=2mm},
    boxed title style={arc=2pt,boxrule=0pt}
  }
}



% ── Cabeçalho / Rodapé ──
\pagestyle{fancy}
\fancyhf{}
\renewcommand{\headrulewidth}{0.4pt}
\renewcommand{\headrule}{\hbox to\headwidth{\color{uefsdark}\leaders\hrule height \headrulewidth\hfill}}
\fancyhead[L]{\small\color{uefsgray}UNIVASF\,--\,Geotecnologias e SIG}
\fancyhead[R]{\small\color{uefsgray}Plano de Ensino 2026.1}
\fancyfoot[C]{\small\color{uefsgray}Página \thepage\ de \pageref{LastPage}}

\hypersetup{colorlinks=true,linkcolor=uefsblue,urlcolor=uefsblue}


\begin{document}

\thispagestyle{empty}




\begin{center}
\color{uefsdark}\rule{\textwidth}{1.2pt}\\[0.30cm]
{\Large\bfseries UNIVERSIDADE FEDERAL DO VALE DO SÃO FRANCISCO}\\[0.15cm]
{\normalsize Colegiado de Geografia}\\[0.06cm]
{\normalsize }\\[0.30cm]
\color{uefsdark}\rule{\textwidth}{1.2pt}
\end{center}

\vspace{0.6cm}

\begin{center}
{\LARGE\bfseries\color{uefsdark} Geotecnologias e SIG}\\[0.30cm]
{\large\bfseries Plano de Ensino\,--\,2026.1}
\end{center}

\vspace{0.4cm}

% ── Identificação ──
\begin{tcolorbox}[uefsbox]
\renewcommand{\arraystretch}{1.25}
\begin{tabularx}{\textwidth}{@{}l X@{}}
\textbf{Carga Horária:} & 11 encontros \\
\textbf{Horário:} & A definir \\
\textbf{Período:} & 2026.1 \\
\textbf{Modalidade:} & Presencial \\
\textbf{Docente:} & Prof.\,Dr.\,Luiz Diego Vidal Santos \\
\textbf{Contato:} & \href{mailto:ldvsantos@uefs.br}{ldvsantos@uefs.br}\enspace·\enspace\href{https://orcid.org/0000-0001-8659-8557}{ORCID: 0000-0001-8659-8557} \\
\end{tabularx}
\end{tcolorbox}

\vspace{0.3cm}



\begin{tcolorbox}[ementabox]
Fundamentos do geoprocessamento e análise espacial; estruturas de dados; análise de incerteza e acurácia; geotecnologias aplicadas a recursos hídricos e seca; SIG na exploração de recursos minerais; intemperismo, erosão e formação de solos; degradação, desertificação e mudanças climáticas; recursos hídricos e gestão por bacia hidrográfica; monitoramento hidrológico; impactos ambientais da mineração; sensoriamento remoto e mudanças de cobertura vegetal; inteligência artificial e qualidade da pesquisa ambiental (princípios FAIR).
\end{tcolorbox}


\section*{Objetivo Geral}

Capacitar o(a) discente a aplicar geotecnologias (SIG, sensoriamento remoto, modelagem espacial e inteligência artificial) na análise, monitoramento e gestão de recursos naturais, com ênfase em rigor metodológico, validação de resultados e suporte à tomada de decisão.

\section*{Objetivos Específicos}

\begin{enumerate}

\item Dominar os fundamentos do geoprocessamento, incluindo estruturas de dados vetoriais e matriciais.

\item Avaliar incerteza e acurácia em produtos de classificação (matriz de confusão, Kappa, validação espacial).

\item Aplicar geotecnologias ao monitoramento de recursos hídricos, seca e evapotranspiração.

\item Utilizar análise multicriterial (AHP, Fuzzy) em SIG para exploração mineral e tomada de decisão.

\item Analisar processos de intemperismo, erosão, formação e degradação de solos com geotecnologias.

\item Delimitar e caracterizar bacias hidrográficas com análise morfométrica.

\item Aplicar sensoriamento remoto na detecção de mudanças de cobertura vegetal.

\item Utilizar aprendizado de máquina com validação espacialmente robusta e princípios FAIR.

\end{enumerate}

\section*{Habilidades e Competências}

\begin{multicols}{2}

\subsection*{Competências Técnicas}

\begin{itemize}

\item Dominar geoprocessamento e análise espacial em plataformas livres (QGIS, GEE)

\item Avaliar qualidade e incerteza de produtos cartográficos

\item Integrar múltiplas fontes de dados espaciais (satelitais, censitários, hidrológicos)

\end{itemize}

\subsection*{Competências Analíticas}

\begin{itemize}

\item Realizar análises multicriteriais e detecção de mudanças

\item Aplicar métodos de ML com validação espacialmente robusta

\item Interpretar indicadores ambientais e séries temporais

\end{itemize}

\columnbreak

\subsection*{Competências Ambientais}

\begin{itemize}

\item Monitorar recursos hídricos e processos de degradação

\item Analisar impactos da mineração com geoprocessamento

\item Avaliar mudanças de cobertura vegetal e métricas de paisagem

\end{itemize}

\subsection*{Competências Profissionais}

\begin{itemize}

\item Elaborar produtos técnicos reprodutíveis (mapas, relatórios, scripts)

\item Aplicar princípios FAIR e ciência aberta na pesquisa ambiental

\item Subsidiar políticas públicas de gestão territorial

\end{itemize}

\end{multicols}

\section*{Significado do Componente para a Formação Profissional}

O componente curricular \textit{Geotecnologias e SIG} é estruturante para a formação de profissionais de Geografia e áreas afins, por fornecer competências em análise espacial, sensoriamento remoto e modelagem ambiental indispensáveis ao diagnóstico, monitoramento e gestão do território. O domínio de geotecnologias é exigência crescente em licenciamentos, auditorias ambientais, planejamento urbano e regional, gestão de recursos hídricos e minerais.

A ênfase em plataformas livres (QGIS, Google Earth Engine), inteligência artificial aplicada e princípios de ciência aberta (FAIR, reprodutibilidade) prepara o(a) egresso(a) para atuar em órgãos públicos, empresas de consultoria, institutos de pesquisa e organizações ambientais.

\section*{Cronograma}


\begin{small}
\renewcommand{\arraystretch}{1.35}
\begin{tabularx}{\textwidth}{@{}c c X@{}}
\toprule
\textbf{Enc.} & \textbf{Sem.} & \textbf{Tema} \\
\midrule
01 & 1\textsuperscript{a} & Incerteza e Acurácia: matriz de confusão, Kappa, validação espacial. \\
02 & 2\textsuperscript{a} & Fundamentos do Geoprocessamento: estruturas vetorial/matricial, operações espaciais. \\
03 & 3\textsuperscript{a} & Geotecnologias e Recursos Hídricos: evapotranspiração, índices de seca, SR hidrológico. \\
04 & 4\textsuperscript{a} & SIG na Exploração Mineral: análise multicriterial, AHP, tomada de decisão. \\
05 & 5\textsuperscript{a} & Intemperismo, Erosão e Solos: processos pedogenéticos, RUSLE, SiBCS. \\
06 & 6\textsuperscript{a} & Degradação, Desertificação e Clima: indicadores espaciais, monitoramento. \\
07 & 7\textsuperscript{a} & Recursos Hídricos e Bacia Hidrográfica: delimitação, morfometria, Lei 9.433/97. \\
08 & 8\textsuperscript{a} & Monitoramento Hidrológico: telemetria IoT, altimetria radar, modelos. \\
09 & 9\textsuperscript{a} & Geoprocessamento na Mineração: DInSAR, NDVI, licenciamento, auditoria. \\
10 & 10\textsuperscript{a} & SR e Cobertura Vegetal: NDVI/EVI, classificadores RF/SVM/U-Net, BFAST. \\
\rowcolor{uefslight!70!white}
\textbf{11} & \textbf{11\textsuperscript{a}} & \textbf{\textsf{IA e Qualidade}: ML, SHAP/LIME, princípios FAIR, Green AI, ética.} \\
\bottomrule
\end{tabularx}
\end{small}

\begin{tcolorbox}[uefsbox]
\textbf{Estrutura geral:}\enspace Enc.\,01--02 → Fundamentos\enspace|\enspace Enc.\,03--06 → Recursos Naturais e Solos\enspace|\enspace Enc.\,07--09 → Hidrologia e Mineração\enspace|\enspace Enc.\,10--11 → SR, IA e Ciência Aberta.
\end{tcolorbox}


\section*{Metodologia}

A disciplina será desenvolvida ao longo de 11 encontros. As atividades combinarão exposição dialogada dos fundamentos teóricos com sessões práticas em ambiente computacional (QGIS, Google Earth Engine e R).

Cada encontro articulará uma dimensão conceitual (fundamentos, pressupostos e critérios) com uma dimensão aplicada (exercícios com dados reais, produção de mapas e relatórios). A disciplina enfatiza o rigor metodológico, a reprodutibilidade dos procedimentos e a interpretação crítica dos resultados.


\begin{tcolorbox}[uefsbox, colframe=uefsaccent, colback=white, title={\bfseries\color{white}\scshape Avaliação}, fonttitle=\normalsize, coltitle=white, colbacktitle=uefsaccent, attach boxed title to top left={yshift=-2mm,xshift=2mm}, boxed title style={arc=2pt,boxrule=0pt}]
\begin{itemize}[leftmargin=1.2em]
\item 1\textsuperscript{a}\,Avaliação (Teórica\,I\,--\,individual): fundamentos, incerteza, recursos hídricos, solos\,\dotfill\,\textbf{25\,\%}
\item 2\textsuperscript{a}\,Avaliação (Prática\,--\,individual/dupla): mapa temático + relatório técnico\,\dotfill\,\textbf{25\,\%}
\item 3\textsuperscript{a}\,Avaliação (Teórica\,II\,--\,individual): bacias, SR, detecção de mudanças, IA\,\dotfill\,\textbf{25\,\%}
\item 4\textsuperscript{a}\,Avaliação (Projeto Final\,--\,individual/grupo): análise integrada reprodutível\,\dotfill\,\textbf{25\,\%}
\end{itemize}
\vspace{4pt}
{\footnotesize Cada avaliação vale até 10,0 pontos. Média final = média aritmética simples.}
\end{tcolorbox}



\section*{Referências}

\subsection*{Básica}

\begin{small}
\hangindent=1.5em\hangafter=1
CÂMARA, G. \textit{et al.} \textbf{Anatomia de Sistemas de Informação Geográfica}. São José dos Campos: INPE, 2004.\par\medskip

\hangindent=1.5em\hangafter=1
NOVO, E.\,M.\,L.\,M. \textbf{Sensoriamento remoto}: princípios e aplicações. 4.\,ed. São Paulo: Blucher, 2008.\par\medskip

\hangindent=1.5em\hangafter=1
ROSA, R.; BRITO, J.\,L.\,S. \textbf{Introdução ao geoprocessamento}. Uberlândia: UFU, 1996.\par\medskip

\hangindent=1.5em\hangafter=1
TUCCI, C.\,E.\,M. \textbf{Hidrologia}: ciência e aplicação. 4.\,ed. Porto Alegre: ABRH/UFRGS, 2002.\par\medskip

\hangindent=1.5em\hangafter=1
ALLEN, R.\,G. \textit{et al.} \textbf{Crop evapotranspiration}. FAO Irrigation and Drainage Paper 56. Rome: FAO, 1998.\par\medskip
\end{small}

\subsection*{Complementar}

\begin{small}
\hangindent=1.5em\hangafter=1
EMBRAPA. \textbf{Sistema Brasileiro de Classificação de Solos}. 5.\,ed. Brasília: Embrapa, 2018.\par\medskip

\hangindent=1.5em\hangafter=1
BRASIL. \textbf{Lei nº 9.433, de 8 de janeiro de 1997}. Política Nacional de Recursos Hídricos. Brasília, 1997.\par\medskip

\hangindent=1.5em\hangafter=1
REICHSTEIN, M. \textit{et al.} Deep learning and process understanding for data-driven Earth system science. \textit{Nature}, v.\,566, p.\,195--204, 2019.\par\medskip

\hangindent=1.5em\hangafter=1
WILKINSON, M.\,D. \textit{et al.} The FAIR Guiding Principles. \textit{Scientific Data}, v.\,3, 160018, 2016.\par\medskip

\hangindent=1.5em\hangafter=1
SCHWARTZ, R. \textit{et al.} Green AI. \textit{Communications of the ACM}, v.\,63, n.\,12, p.\,54--63, 2020.\par\medskip
\end{small}



\vfill

\begin{center}
\noindent\rule{6cm}{0.4pt}\\[6pt]
\textbf{Prof.\,Dr.\,Luiz Diego Vidal Santos}\\
Docente Responsável\\[4pt]
{\small Feira de Santana\,--\,BA, fevereiro de 2026.}
\end{center}

\end{document}
